\section{Introduction}
\label{sec:intro}

ERA5~\cite{hersbach2020} exceeds 400\,TB in float32, making storage and data
movement practical constraints. CRA5~\cite{han2024,cra5journal} uses a transformer
variational autoencoder to reduce this archive to about 0.85\,TB. Its published
evaluation describes reconstruction error by variable, pressure level, spatial
location, and spectrum. A separate question is how much error remains predictable
from the decoded fields themselves.

We attach a 1.65\,M-parameter U-Net to CRA5's decoded output.\footnote{Code and reproducibility: \url{https://github.com/tj567356/CRA5-Decoder-Correction}.}
On the 2018 CONUS test year, mean normalised RMSE (nRMSE) decreases by 25.8\% over the 19
FourCastNet variables directly represented by CRA5, without changing the
per-frame bitstream. Including derived total column water vapour gives 27.8\%.
We call this reduction \emph{model-recoverable error}: it depends on the chosen
network and training protocol and is not an information-theoretic bound.
On the complete 2022 year outside CRA5's published training window, the same
frozen corrector reduces direct-channel nRMSE by 26.4\% without fine-tuning.

The reduction is uneven: 32.2\% for near-surface fields versus 3.5\% aloft.
Their 9.1-fold ratio compares correction gains, not compression ratios. We apply
the same correction architecture, with separately fitted weights, to classical
transform codecs. A single-stage MPCA-only codec matched to CRA5's mean bits per
value gives a near/upper gain ratio of 1.7; three MPCA+PCA settings at higher
compression give 1.32--1.35. These controls help assess whether the asymmetry is
specific to the tested codec and corrector combination.

Terrain is a plausible contributor: small-scale topography may be difficult
to represent on a coarse latent grid, but error magnitude and correction gain
are different quantities. Using actual ERA5 orography and land--sea masks, we
find larger baseline errors and larger gains on rougher western land: $G$ rises
from 18.1\% to 36.9\% between roughness endpoint deciles. The association
persists within several elevation bands but does not isolate a terrain mechanism.
Our earlier proxy-elevation analysis remains a sensitivity check.

We contribute (i) a CRA5/MPCA comparison using measured bitstreams at
matched mean rate, distinguishing global and regional coding; (ii) vertical
recoverability with frozen 2022 testing and corrector controls; and (iii)
ERA5-orography roughness associations with elevation and geographic checks.
Checkpoints are selected on 2017 data; temporal block-bootstrap intervals
condition on each fitted model and test year.

